\maketitle

\begin{abstract}
We prove the conjecture that the irrationality exponent of \(\pi\) is \(2\).
As a consequence, the classical Flint--Hills series
\(\sum_{n\ge1}1/(n^3\sin^2 n)\) converges, with angles in radians.
\end{abstract}

\section{Introduction}\label{sec:introduction}

For an irrational real number \(x\), its \emph{irrationality exponent}
(also called its irrationality measure) is
\[
\mu(x)=\sup\left\{\nu>0:
0<\left|x-\frac pq\right|<q^{-\nu}
\text{ for infinitely many coprime }p,q\in\ZZ,\ q\ge2\right\}.
\]
The pigeonhole principle gives \(\mu(x)\ge2\). An upper bound limits
how frequently rational numbers can approximate \(x\) exceptionally
closely; irrationality alone supplies no finite upper bound. The expected
value for the usual positive circle constant is \(2\). Waldschmidt
records the corresponding epsilon-dependent approximation conjecture
in \cite[p.~265]{waldschmidt}. Our main result proves it.

\begin{theorem}\label{thm:main}
The irrationality exponent of \(\pi\) is \(2\). More precisely, for every
real \(\nu>2\), there is an integer \(Q(\nu)\) such that
\[
\left|\pi-\frac pq\right|\ge q^{-\nu}
\qquad(p\in\ZZ,\ q\in\ZZ,\ q\ge Q(\nu)).
\]
\end{theorem}

The inequality applies to all integer numerators and denominators,
whether or not the fraction is reduced. The threshold may depend on
\(\nu\), and the argument does not give it effectively. In particular,
an exponent of \(2\) is different from a uniform positive lower bound
of the form \(c/q^2\), or equivalently from bounded continued-fraction
partial quotients.

The related Flint--Hills problem, popularized by Pickover and discussed
by Alekseyev \cite{alekseyev}, asks whether
\[
\sum_{n\ge1}\frac{1}{n^3\sin^2 n}
\]
converges, with angles in radians. Integers very close to multiples of
\(\pi\) make individual denominators small. A rational-approximation
bound controls how often such small values can occur, not just their
individual size.

\begin{corollary}\label{thm:flinthills}
The classical Flint--Hills series converges.
\end{corollary}

Alekseyev showed that convergence requires
\(\mu(\pi)\le5/2\) \cite[Corollary~4]{alekseyev};
